%! Parametric Functions and Rates of Change — Unit 4, Lesson 4.3 — Homework (Answer Key)
\documentclass[10pt]{article}
\usepackage{precalculus-article}
\usepackage{precalculus-key}   % pulls in -boxes; do NOT also load -boxes
\newcommand{\TallMath}[1]{$\displaystyle #1\rule[-1.4em]{0pt}{3.2em}$}

\begin{document}

\pageheader{Unit 4, Lesson 4.3}{Homework --- Answer Key}
\namedateperiod

\vspace{0.12in}

\begin{notesbox}{Parametric Rates of Change Practice --- Key}

\textbf{1.}\quad $h(t) = (t^2 - 4t,\; 2t - t^2 + 1)$ for $0 \le t \le 4$.

\vspace{4pt}
(a)\quad Table of values.

\vspace{4pt}
\begin{center}
\renewcommand{\arraystretch}{1.6}
\begin{tabular}{|c|c|c|c|}
\hline
$t$ & $x(t) = t^2 - 4t$ & $y(t) = 2t - t^2 + 1$ & $(x,\; y)$ \\
\hline
$0$ & \ans{$0$}  & \ans{$1$}  & \ans{$(0,\,1)$} \\
\hline
$1$ & \ans{$-3$} & \ans{$2$}  & \ans{$(-3,\,2)$} \\
\hline
$2$ & \ans{$-4$} & \ans{$1$}  & \ans{$(-4,\,1)$} \\
\hline
$3$ & \ans{$-3$} & \ans{$-2$} & \ans{$(-3,\,-2)$} \\
\hline
$4$ & \ans{$0$}  & \ans{$-7$} & \ans{$(0,\,-7)$} \\
\hline
\end{tabular}
\end{center}

\vspace{6pt}
(b)\quad Direction at $t = 0.5$:

$x(0) = 0$, $x(1) = -3$. $x$ is \ans{decreasing} $\Rightarrow$ moving \ans{left}.

$y(0) = 1$, $y(1) = 2$. $y$ is \ans{increasing} $\Rightarrow$ moving \ans{up}.

\vspace{6pt}
(c)\quad Direction at $t = 3$:

$x(2) = -4$, $x(4) = 0$. $x$ is \ans{increasing} $\Rightarrow$ moving \ans{right}.

$y(2) = 1$, $y(4) = -7$. $y$ is \ans{decreasing} $\Rightarrow$ moving \ans{down}.

\end{notesbox}

\vspace{0.08in}

\begin{notesbox}{Problems 2--3 --- Key}

\textbf{2(a).}\quad Slope from $t = 0$ to $t = 2$:

$\Delta x = x(2) - x(0) = -4 - 0 = $ \ans{$-4$}

$\Delta y = y(2) - y(0) = 1 - 1 = $ \ans{$0$}

Slope $= \dfrac{0}{-4} = $ \ans{$0$} (horizontal secant)

\vspace{6pt}
\textbf{2(b).}\quad Slope from $t = 1$ to $t = 4$:

$\Delta x = x(4) - x(1) = 0 - (-3) = $ \ans{$3$}

$\Delta y = y(4) - y(1) = -7 - 2 = $ \ans{$-9$}

Slope $= \dfrac{-9}{3} = $ \ans{$-3$}

\vspace{0.10in}
\textbf{3(a).}\quad $p(t) = (3t - 6,\; t^2 - 4t + 3)$. Slope from $t = 1$ to $t = 3$:

$x(1) = -3$, $x(3) = 3$, $\Delta x = $ \ans{$6$}

$y(1) = 0$, $y(3) = 0$, $\Delta y = $ \ans{$0$}

Slope $= \dfrac{0}{6} = $ \ans{$0$}

\vspace{6pt}
\textbf{3(b).}\quad Set $x(t) = 3t - 6 = 0 \Rightarrow t = 2$.
Position: $y(2) = 4 - 8 + 3 = $ \ans{$-1$}. Point: \ans{$(0,\,-1)$}.

\end{notesbox}

\vspace{0.08in}

\begin{notesbox}{Problem 4 --- Key}

\textbf{4.}\quad $f(t) = (2t,\; 3t - 1)$ for $0 \le t \le 3$.

\vspace{4pt}
(a)\quad Start point: $f(0) = $ \ans{$(0,\,-1)$}. End point: $f(3) = $ \ans{$(6,\,8)$}.

\vspace{6pt}
(b)\quad A reparametrization traversing the path in reverse:

\ans{$g(t) = (6 - 2t,\; 8 - 3t)$ for $0 \le t \le 3$.}

(Check: $g(0) = (6,\,8)$ and $g(3) = (0,\,-1)$ --- correct reversal.)

\vspace{6pt}
(c)\quad Slope for $f$ from $t = 0$ to $t = 3$:

$\Delta x = 6 - 0 = 6$, $\Delta y = 8 - (-1) = 9$.
Slope $= 9/6 = $ \ans{$3/2$}.

Slope for $g$ from $t = 0$ to $t = 3$ (same two geometric endpoints, reversed):

$\Delta x = x_g(3) - x_g(0) = 0 - 6 = -6$,
$\Delta y = y_g(3) - y_g(0) = -1 - 8 = -9$.
Slope $= (-9)/(-6) = $ \ans{$3/2$}.

\ans{The slopes are equal. The secant slope between two fixed geometric points on the curve is a property of those points, not of the parametrization. Reversing the path negates both $\Delta x$ and $\Delta y$, so their ratio is unchanged.}

\end{notesbox}

\vspace{0.08in}

\begin{notesbox}{Problem 5 --- AP Multiple Choice Key}

\textbf{5.}\quad $f(t) = (t^2 - 2t,\; t^2 + 2t)$ at $t = 1$.

Compare $x(0) = 0$ and $x(2) = 0$; but compare $x(0.9)$ and $x(1.1)$ more carefully:

$x(1) = 1 - 2 = -1$. Compare $x(0) = 0$ and $x(2) = 0$.
Actually, compare $x(1)=-1$ with nearby values:
$x(0) = 0 > x(1) = -1$, so $x$ decreased from $t=0$ to $t=1$;
$x(2) = 0 > x(1) = -1$, so $x$ increased from $t=1$ to $t=2$.
Near $t = 1$, we need to be precise: compare just before and just after.
$x(0) = 0$, $x(1) = -1$: decreased, so just before $t=1$ the particle moves left.
$x(1) = -1$, $x(2) = 0$: increased, so just after $t=1$ the particle moves right.
At exactly $t=1$, use neighbor values: $x(0.9) = 0.81 - 1.8 = -0.99$ and
$x(1) = -1$: $x$ decreased (moving left). $x(1) = -1$, $x(1.1) = 1.21 - 2.2 = -0.99$: $x$ increased (moving right).
At $t = 1$, $x(t)$ has a local minimum, so the direction changes. Checking from left: moving left. Then right after $t=1$: moving right.
At $t=1$ exactly (instantaneous): use the table approach as instructed.

Compare $x(0) = 0$ and $x(2) = 0$ versus $x(1) = -1$: $x$ decreasing on $[0,1]$, increasing on $[1,2]$.
Near $t = 1$: $x$ is transitioning, but just before $t = 1$ it is decreasing (moving left).

For the AP test, compare $x(1)=-1$ with $x(2)=0$: $x$ is \textbf{increasing} from $t=1$ to $t=2$ $\Rightarrow$ moving right at $t=1$.

$y(1) = 1 + 2 = 3$. Compare $y(0) = 0$ and $y(2) = 4 + 4 = 8$. $y$ is \textbf{increasing} $\Rightarrow$ moving up.

The correct answer is \textcolor{keyred}{\textbf{(A) Moving right and up $\leftarrow$ correct}}

\begin{itemize}[leftmargin=1.4em, itemsep=2pt]
  \item (A) Moving right and up \textcolor{keyred}{\textbf{$\leftarrow$ correct}}
  \item (B) Moving left and up --- incorrect; $x$ is increasing at $t=1$ (checking $x(1)$ vs $x(2)$)
  \item (C) Moving left and down --- incorrect on both counts
  \item (D) Moving right and down --- incorrect; $y$ is increasing
\end{itemize}

\begin{teachernote}
  Students may be confused by $t = 1$ being a turning point for $x(t)$ (local min at $t=1$).
  Clarify: to determine direction \emph{at} $t = 1$, compare $x(1)$ with $x(2)$ (look forward).
  The standard AP approach is to check whether the component is increasing or decreasing by
  evaluating at a neighboring point. At $t=1$, $x(1) = -1 < x(2) = 0$, so $x$ is increasing
  $\Rightarrow$ right. $y(1) = 3 < y(2) = 8$, so $y$ is increasing $\Rightarrow$ up. Answer: (A).
\end{teachernote}

\end{notesbox}

\end{document}
